problem:  
Let \( G \) be a compact simple Lie group, \( H \subset G \) a closed connected reductive subgroup, and \( G/H \) a compact homogeneous space admitting \( G \)-invariant Einstein metrics. Write \( \mathfrak{g} = \mathfrak{h} \oplus \mathfrak{m} \) with \(\mathfrak{m} = \bigoplus_{i=1}^r \mathfrak{m}_i\) the irreducible decomposition of the isotropy representation.  

For each \( G \)-invariant metric \(g\) determined by positive weights \( x_1, \dots, x_r \) on the \( \mathfrak{m}_i\), let \( \mathrm{Ric}_g:\mathfrak{m}\to\mathfrak{m} \) denote the Ricci operator, and let  
\[
\sigma(g) = \text{spec}\big(\widehat{\mathrm{Ric}}_g\big)
\]
denote its *normalized spectrum*, where \( \widehat{\mathrm{Ric}}_g = (\mathrm{scal}(g))^{-1}\mathrm{Ric}_g \) is scaled so that scalar curvature equals one.  

**Problem:**  
Investigate whether, for multiplicity-free isotropy representation (i.e. all \(\mathfrak{m}_i\) pairwise nonisomorphic as \(H\)-modules), the map  
\[
g \;\mapsto\; \sigma(g)
\]
is injective on the moduli of \( G \)-invariant Einstein metrics up to isometry and scaling.  
Equivalently: does the normalized Ricci spectrum completely determine a homogeneous Einstein metric in the multiplicity-free reductive setting?  
If not, characterize the explicit algebraic conditions on the structure constants of the isotropy decomposition allowing distinct invariant Einstein metrics to share the same normalized Ricci spectrum.

Why is it a "good" problem:  
This problem isolates a subtle rigidity question linking homogeneous Einstein geometry, curvature operators, and representation theory of isotropy subgroups. It is deceptively simple in form—asking whether the normalized Ricci curvature spectrum uniquely determines a \(G\)-invariant Einstein metric—but deeply nontrivial in content. Invariant Einstein metrics correspond to algebraic critical points of a highly nonlinear polynomial system in the weights \(x_i\). The Ricci spectrum captures only eigenvalue data of this system. Determining whether spectral data suffice to recover geometric equivalence probes the boundary between algebraic rigidity and spectral ambiguity in reductive geometry.  
Solving—or even partially analyzing—this question would bridge classification theory of homogeneous Einstein spaces (in the spirit of Wang–Ziller and Böhm–Kerr) with spectral geometry, possibly revealing new moduli or hidden symmetries. It is precisely in the multiplicity-free case that conceptual clarity arises yet known classification results remain incomplete, making the problem both precise, logically sound, and genuinely open.